\documentclass[11pt,a4paper]{article}
\usepackage[T1]{fontenc}
\usepackage{lmodern}
\usepackage[a4paper,margin=24mm]{geometry}
\usepackage{amsmath,amssymb,mathtools}
\usepackage{xcolor,fancyhdr,hyperref}
\definecolor{ink}{HTML}{18354A}
\hypersetup{colorlinks=true,linkcolor=ink,urlcolor=ink,
  pdftitle={Direct equivalence of the two C4 rules},
  pdfsubject={Syntactic derivations for Figure1 and Figure1-MS}}
\pagestyle{fancy}
\fancyhf{}
\fancyhead[L]{\small\textcolor{ink}{Figure1 $\longleftrightarrow$ Figure1-MS}}
\fancyhead[R]{\small Direct C4 derivations}
\fancyfoot[C]{\small\thepage}
\setlength{\headheight}{14pt}
\setlength{\parindent}{0pt}
\setlength{\parskip}{6pt}
\setlength{\abovedisplayskip}{8pt}
\setlength{\belowdisplayskip}{8pt}
\newcommand{\F}{\mathbb F_p}
\newcommand{\eqw}{\mathrel{\simeq}}
\newcommand{\CC}{\mathcal C}
\newcommand{\rulebox}[2]{\begin{center}\fcolorbox{ink}{ink!4}{%
  \parbox{0.9\linewidth}{\centering\textbf{#1}\quad $\displaystyle #2$}}\end{center}}
\newcommand{\heading}[1]{\section*{\textcolor{ink}{#1}}}
\begin{document}

{\LARGE\bfseries\textcolor{ink}{Direct equivalence of the two C4 rules}}\par
\smallskip
{\large Figure1 and Figure1-MS, with scalars retained}

Let $p$ be an odd prime and $g\in\F^{\times}$ a primitive root. Write
\[
 M=M_g,\qquad u=g^{-1},\qquad d=u^2,\qquad b=g^2,
 \qquad c=\frac{1-g}{2}\,d,\qquad e=\frac{g^2-g}{2}.
\]
Every displayed exponent of $S$ or $Z$ is an element of $\F$; the word
$S^a$ means the natural-number power at the representative of $a$ in
$\{0,\ldots,p-1\}$. Fractions and inverses are computed in $\F$.
Juxtaposition is word concatenation, with associativity and units suppressed.
The symbol $\eqw$ denotes derivability from the specified axioms, and $1$
denotes the empty word. These are exact equations, not equations modulo scalars.

\rulebox{Original C4 (Figure1)}{MS\eqw Z^cS^dM}
\rulebox{Replacement C4 (Figure1-MS)}{SM\eqw MS^bZ^e}

\heading{1. Consequences of the shared axioms}
Both presentations share C0--C3, C5--C15 and the structural circuit rules.
The one-wire support needed below follows from C0--C3, C5 and scalar centrality:
\[
 S^p\eqw Z^p\eqw1,\qquad ZS\eqw SZ,\qquad
 J^2\eqw1,\qquad JM\eqw MJ,\qquad J=H^2.
 \tag{1}
\]
Here is the derivation of this support. C3 gives
$M_g^k\eqw M_{g^k}$, $M_1\eqw1$ and $M_g^{p-1}\eqw1$.
Taking discrete logarithms of nonzero field elements and reducing exponents
modulo $p-1$ gives
\[
 M_xM_y\eqw M_{xy},\qquad M_{-1}^2\eqw1,\qquad
 M_{-1}M\eqw MM_{-1}.
\]
Set $\lambda=(\nu^p)^{(p-1)/2}$, with $\nu$ the central scalar generator.
C0 gives $\nu^{2p}\eqw1$, hence $\lambda^2\eqw1$.
C2 says $J\eqw M_{-1}\lambda$, so $J^2\eqw1$ and $JM\eqw MJ$.

Put $T=JSJ$. C5 says $ST\eqw TS$. Since
$Z=JSJS^{-1}=TS^{-1}$, it follows that $ZS\eqw SZ$.
Also $T^p\eqw JS^pJ\eqw1$ by C1 and $J^2\eqw1$; thus
$Z^p\eqw T^pS^{-p}\eqw1$. This proves (1) without either C4.

\newpage
\heading{2. Deriving the action on $Z$ from each C4}
Define $\CC(w)=JwJ$. By (1), this is an involutive word automorphism:
\[
 \CC(xy)\eqw\CC(x)\CC(y),\quad
 \CC^2(w)\eqw w,\quad \CC(M)\eqw M.
\]
The definition of $Z$ gives $\CC(S)\eqw ZS$.
Applying $\CC$ again yields
$S\eqw\CC(Z)ZS$. Cancelling the final $S$ gives
$\CC(Z)Z\eqw1$, so
\[
 \CC(Z)\eqw Z^{-1},\qquad
 \CC(Z^aS^h)\eqw Z^{h-a}S^h.
 \tag{2}
\]
All cancellations here and below use the finite-order inverses from (1).

\subsection*{Starting with the original C4}
Assume $MS\eqw Z^cS^dM$. Apply (2):
\begin{align*}
 MZS
 &\eqw \CC(MS)
 \eqw \CC(Z^cS^dM)
 \eqw Z^{d-c}S^dM\\
 &\eqw Z^{d-2c}(Z^cS^dM)
 \eqw Z^{d-2c}MS.
\end{align*}
Cancel $S$. Since
\[
 c=\frac{d-u}{2},\qquad d-2c=u,
\]
we obtain $MZ\eqw Z^uM$. Iteration gives
$MZ^x\eqw Z^{ux}M$. Taking $x=gk$ and using $ug=1$ gives
\[
 \boxed{Z^kM\eqw MZ^{gk}}.
 \tag{3}
\]

\subsection*{Starting with the replacement C4}
Assume $SM\eqw MS^bZ^e\eqw MZ^eS^b$.
Apply (2), then substitute the assumed rule:
\begin{align*}
 ZSM&\eqw MZ^{b-e}S^b,\\
 (ZM)(Z^eS^b)
 &\eqw ZSM
 \eqw MZ^{b-e}S^b
 \eqw (MZ^{b-2e})(Z^eS^b).
\end{align*}
Cancel $Z^eS^b$. As $b-2e=g$, this gives $ZM\eqw MZ^g$.
Iteration again yields (3). In either theory, taking $k=ux$ in (3) also gives
\[
 \boxed{MZ^x\eqw Z^{ux}M}.
 \tag{4}
\]
Thus each theory supplies (3)--(4) from its own C4; the argument is not circular.

\newpage
\heading{3. Original C4 implies the replacement C4}
Assume only the original C4 and the shared axioms:
\[
 MS\eqw Z^cS^dM.
\]
Use (4), derived from this assumption on the preceding page.

\subsection*{Step 1: iterate the assumed rule $b=g^2$ times}
Choose the natural representative of $b$ for the induction. Because $Z$ and
$S$ commute and have order $p$, their powers add and multiply in $\F$.
Hence
\begin{align*}
 MS^b
 &\eqw (Z^cS^d)^bM
 &&\text{iterate the original C4}\\
 &\eqw Z^{cb}S^{db}M
 &&\text{$ZS\eqw SZ$ and power laws}\\
 &\eqw Z^{cb}SM
 &&\text{$db=u^2g^2=1$.}
\end{align*}

\subsection*{Step 2: move the remaining $Z$-power through $M$}
\begin{align*}
 MS^bZ^e
 &\eqw Z^{cb}SMZ^e
 &&\text{Step 1}\\
 &\eqw Z^{cb}SZ^{eu}M
 &&\text{(4) with $x=e$}\\
 &\eqw Z^{cb+eu}SM
 &&\text{$ZS\eqw SZ$}\\
 &\eqw SM
 &&\text{$cb+eu=0$.}
\end{align*}
The last field identity is
\[
 cb+eu
 =\frac{1-g}{2}\,u^2g^2+\frac{g^2-g}{2}\,u
 =\frac{1-g}{2}+\frac{g-1}{2}=0.
\]
Reverse the resulting equation to obtain exactly the replacement C4:
\rulebox{Derived in Figure1}{SM\eqw MS^{g^2}Z^{(g^2-g)/2}}

\subsection*{Formal counterpart}
In \texttt{FF-Iso.agda}, the iteration and calculation above are
\texttt{Conversion.MS}\,$\Rightarrow$\,\texttt{MR}.
The theorem \texttt{Figure1.c4'} reverses its conclusion to match the
orientation of the Figure1-MS axiom.

\newpage
\heading{4. Replacement C4 implies the original C4}
Assume only the replacement C4 and the shared axioms:
\[
 SM\eqw MS^bZ^e.
\]
Use (4), now derived from this assumption. Set $f=ed$.

\subsection*{Step 1: iterate the assumed rule $d=g^{-2}$ times}
Again choose the natural representative for the induction. Then
\begin{align*}
 S^dM
 &\eqw M(S^bZ^e)^d
 &&\text{iterate the replacement C4}\\
 &\eqw MS^{bd}Z^{ed}
 &&\text{$ZS\eqw SZ$ and power laws}\\
 &\eqw MSZ^f
 &&\text{$bd=g^2u^2=1$.}
\end{align*}

\subsection*{Step 2: remove the extra $Z$-power}
\begin{align*}
 MS
 &\eqw (MSZ^f)Z^{-f}
 &&\text{insert $Z^fZ^{-f}\eqw1$}\\
 &\eqw S^dMZ^{-f}
 &&\text{Step 1}\\
 &\eqw S^dZ^{-fu}M
 &&\text{(4) with $x=-f$}\\
 &\eqw Z^cS^dM
 &&\text{$ZS\eqw SZ$ and $-fu=c$.}
\end{align*}
The last field identity is
\[
 -fu=-edu
 =-\frac{g^2-g}{2}\,u^3
 =\frac{u^2-u}{2}
 =\frac{1-g}{2}\,u^2=c.
\]
This is exactly the original C4:
\rulebox{Derived in Figure1-MS}{MS\eqw
   Z^{(1-g)g^{-2}/2}S^{g^{-2}}M}

\subsection*{Formal counterpart and equivalence}
This direction is \texttt{Conversion.MR}\,$\Rightarrow$\,\texttt{MS} and
\texttt{Figure1-MS.c4} in \texttt{FF-Iso.agda}.
Every other axiom is carried over unchanged, and derivations are transported
recursively through the congruence and structural rules. Therefore the two
presentations generate the same equivalence relation on the same circuits.
Their presented groups are isomorphic by the identity on words.

\smallskip
\small
\textbf{Formal source:}
\nolinkurl{Examples/Groups/Clifford+MinusOne/Qupit/Figure1-MS/FF-Iso.agda}.
The Agda file proves both directions syntactically, without importing Clifford
semantics or the existing exact-group isomorphisms.

\end{document}
